\documentclass[11pt]{article}
\usepackage[a4paper,margin=1in]{geometry}
\usepackage{fontspec}
\usepackage[boldfont=false]{xeCJK}
\setCJKmainfont{FandolSong-Regular.otf}
\usepackage{amsmath}
\usepackage{amssymb}
\usepackage{array}
\usepackage{booktabs}
\usepackage{graphicx}
\usepackage{adjustbox}
\usepackage{enumitem}
\setlist[itemize]{topsep=2pt,partopsep=0pt,parsep=0pt,itemsep=2pt,leftmargin=1.4em}
\setlist[enumerate]{topsep=2pt,partopsep=0pt,parsep=0pt,itemsep=2pt,leftmargin=1.6em}
\usepackage{parskip}
\usepackage{fvextra}
\fvset{breaklines=true,breakanywhere=true,fontsize=\small}
\setlength{\parindent}{0pt}

\begin{document}

\begin{center}
  {\LARGE\bfseries Polymerisation}\\[0.35em]
  {\large A Level Chemistry}
\end{center}
\vspace{0.6em}

\section*{Addition polymerisation}

In \textbf{addition polymerisation} 加成聚合, many small molecules join into one very long chain, with no other product made.

Each small molecule is a \textbf{monomer} 单体. It must be unsaturated — it has a C=C double bond. The double bond opens up so that the monomers can link together. The long chain that forms is the \textbf{polymer} 聚合物.

\subsection*{Repeat units}

The \textbf{repeat unit} 重复单元 is the small part that is copied again and again along the chain. To find it, take the monomer, change the C=C to a single C–C, and draw bonds going out at each end.

\begin{itemize}
\item \textbf{poly(ethene)} 聚乙烯 is made from ethene:

\end{itemize}

\[
n\,\text{CH}_2{=}\text{CH}_2 \rightarrow -(\text{CH}_2{-}\text{CH}_2)_n-
\]

\begin{itemize}
\item \textbf{poly(chloroethene)} 聚氯乙烯 (PVC) is made from chloroethene:

\end{itemize}

\[
n\,\text{CH}_2{=}\text{CHCl} \rightarrow -(\text{CH}_2{-}\text{CHCl})_n-
\]

\subsection*{Finding the monomer}

To go the other way, look at one repeat unit of the polymer, and put the C=C double bond back in. That gives you the monomer.

\section*{The problem of disposal}

Poly(alkene)s are very hard to get rid of:

\begin{itemize}
\item they are \textbf{non-biodegradable} 不可生物降解 — microbes cannot break them down, so they stay in the ground for a very long time.

\item their \textbf{combustion} 燃烧 (burning) can release harmful gases. For example, burning PVC gives off toxic hydrogen chloride.

\end{itemize}


\end{document}
